% from ega4-16-fr.tex
\oldpage[IV-4]{5}
Dans ce paragraphe, nous présentons, sous forme globale, quelques notions de
calcul différentiel particulièrement utiles en Géométrie algébrique. Nous passons sous
silence de nombreux développements, classiques en Géométrie différentielle (connexions,
transformations infinitésimales associées à un champ de vecteurs, jets, etc.), bien que ces
notions s'écrivent de façon particulièrement naturelle dans le cadre des schémas. Nous
passons également sous silence ici les phénomènes spéciaux à la caractéristique $p>0$
(dont certains sont étudiés, dans le cadre affine, dans (0, 21)). Pour certains compléments
sur le formalisme différentiel dans les préschémas, le lecteur pourra consulter les
exposés II et VII de [42], ainsi que les chapitres ultérieurs de ce Traité.

\subsection{Invariants normaux d'une immersion}
\EGAReaderTarget{ega4.c04.subsec.16-1}{IV.16.1}

\begin{env}[16.1.1]
\label{IV.16.1.1}
Soient $(X,\mathcal{O}_X)$, $(Y,\mathcal{O}_Y)$ deux espaces annelés,
$f=(\psi,\theta):Y\to X$ un morphisme d'espaces annelés (0\textsubscript{I}, 4.1.1)
tel que l'homomorphisme
\[
  \theta^\#:\psi^*(\mathcal{O}_X)\longrightarrow\mathcal{O}_Y
\]
soit surjectif, de sorte que $\mathcal{O}_Y$ s'identifie à un faisceau d'anneaux quotient
$\psi^*(\mathcal{O}_X)/\mathscr{I}_f$. On peut alors munir
$\psi^*(\mathcal{O}_X)$ de la filtration $\mathscr{I}_f$-préadique.
\end{env}

\begin{definition}[16.1.2]
\EGAReaderTarget{ega4.c04.def.16-1-2}{IV.16.1.2}
Le faisceau d'anneaux $\mathcal{O}_Y$-augmenté
$\psi^*(\mathcal{O}_X)/\mathscr{I}_f^{n+1}$ est appelé le $n$-ème invariant normal
de $f$; l'espace annelé
$(Y,\psi^*(\mathcal{O}_X)/\mathscr{I}_f^{n+1})$ est appelé le $n$-ème voisinage
infinitésimal de $Y$ pour le morphisme $f$, et noté $Y_f^{(n)}$ ou simplement
$Y^{(n)}$. Le faisceau d'anneaux gradués associé au faisceau d'anneaux filtrés
$\psi^*(\mathcal{O}_X)$
\[
  \mathscr{G}r_\bullet(f)=\bigoplus_{n\geq 0}
  (\mathscr{I}_f^n/\mathscr{I}_f^{n+1})
  \tag{16.1.2.1}
  \label{IV.16.1.2.1}
\]
est appelé le faisceau d'anneaux gradués associé à $f$. Le faisceau
$\mathscr{G}r_1(f)=\mathscr{I}_f/\mathscr{I}_f^2$ est appelé le faisceau conormal
de $f$ (que l'on note aussi $\mathscr{N}_{Y/X}$ s'il n'en résulte pas de confusion).
\end{definition}

Il est clair que les
$\mathcal{O}_{Y^{(n)}}=\psi^*(\mathcal{O}_X)/\mathscr{I}_f^{n+1}$
(qu'on note aussi $\mathcal{O}_{Y_f^{(n)}}$) forment un
